\documentclass[10pt,a4paper]{article}
\usepackage[utf8]{inputenc}
\usepackage[T1]{fontenc}
\usepackage[margin=18mm]{geometry}
\usepackage{longtable}
\usepackage{hyperref}
\usepackage{enumitem}
\setlist[itemize]{noitemsep, topsep=2pt}
\setlength{\parindent}{0pt}
\begin{document}

\section*{Non-Verbal \& Multimodal Discourse Analysis Worksheet}

\textit{Qualitative analysis of pauses, trailing sentences, prosody/intonation, vocalisations (coughs, sighs, scoffs), and kinesics (hand gestures, facial expressions, posture) in political news discourse}\par\smallskip

\bigskip\hrule\bigskip

\subsection*{Guidelines: Five non-verbal \& paralinguistic features to look for in every event}

\begin{longtable}{|p{0.23\textwidth}|p{0.23\textwidth}|p{0.23\textwidth}|p{0.23\textwidth}|}
\hline
Feature & What to look for & Illustrative example (from sample) & Questions to answer \\ \hline \endhead
\textbf{1. Pauses, Silences \& Abrupt Stops} & Micro-pauses \texttt{(.)}, timed silences \texttt{(1.5)}, pregnant/dramatic pauses before a loaded word, hesitant stumbles, abrupt self-interruptions or cut-offs (\texttt{---}), and forced silences (muting a panelist's mic or talking over them). & Anchor pauses \texttt{(1.5)} before saying \textit{"'loota'"} for dramatic weight; Addl. DCP cuts off mid-sentence (\textit{"wo toh--- abhi investigation chal raha hai"}) when asked if police struck first; panelist's mic faded during counter-question. & Where does the pause or stop occur? Is it deliberate/rhetorical (for suspense, emphasis, or gravitas) or defensive/evasive (cognitive load, self-repair, avoiding blame)? Who controls turn-taking and silence? \\ \hline
\textbf{2. Trailing Sentences \& Unfinished Utterances} & Clauses that trail off (\texttt{...}), drop in volume at the end, or leave the conclusion unspoken (aposiopesis) so the audience supplies the implication without the speaker taking explicit responsibility. & \textit{"Aur phir wahan jo log ghuse the... aap samajh hi sakte hain..."}; \textit{"500 salon ke sangharsh ke baad agar aisa..."} trailing off with a slow head shake. & What exact thought is left unfinished? What presupposition or shared bias is the viewer invited to fill in? Does trailing off act as a safe hedge to imply higher-level complicity or "anti-national" motive without stating it? \\ \hline
\textbf{3. Intonation, Pitch, Volume \& Prosody} & Pitch contour (rising \texttt{$\uparrow$} for incredulity/challenge; falling \texttt{$\downarrow$} for finality), volume shifts (shouting vs. hushed whisper), emphatic word stress, speech rate (rapid-fire urgency vs. slow solemnity), and tonal register (sarcastic, reverential, combative). & Hushed, reverential cadence on \textit{"Prabhu Shri Ram ke Mandir"} vs. sharp, high-pitched interrogation on \textit{"par Akhilesh ji Ram Mandir kyon$\uparrow$ nahin jate?!"}; satirical sing-song inflection on \textit{"priventiv arenjment"}. & Which exact words receive pitch peaks or volume stress? How does vocal tone turn neutral words into irony, mockery, or moral outrage? Does the speaker shift prosodic register depending on which actor is discussed? \\ \hline
\textbf{4. Coughs, Sighs \& Non-Lexical Vocalisations} & Audible sighs (\texttt{[sigh]}), scoffs/snorts (\texttt{[scoff]}), ironic chuckles (\texttt{[chuckle]}), throat-clearing/coughs (\texttt{[cough]}), sharp intakes of breath/gasps (\texttt{[gasp]}), and tongue-clicks (\texttt{[tsk]}). & Heavy \texttt{[sigh]} before reporting Trust oversight lapses (signalling reluctant sorrow rather than hostility); mocking \texttt{[scoff]} or \texttt{[chuckle]} when reading an opponent's denial; \texttt{[cough / throat-clear]} when pivoting away from an awkward fact. & What emotion or epistemic stance does the vocalisation perform (exasperation, moral disgust, disbelief, ridicule, discomfort)? Does a sigh or scoff pre-judge a speaker's claim before any verbal rebuttal is made? \\ \hline
\textbf{5. Hand Gestures, Facial Expressions \& Posture (Kinesics)} & Hand movements (air-quotes, finger-pointing/wagging, desk-thumping, open palms, counting on fingers, dismissive wave), facial cues (raised eyebrows, eye-rolls, smirk, furrowed brow), head nods/shakes, leaning into camera, and visual props (bandaged limbs, documents). & Anchor forms physical air-quotes on \textit{"'loota'"} or \textit{"'peaceful' protest"}; host points finger/pen at lens during an ultimatum; injured DCP holds bandaged hand up to camera; creator raises one eyebrow and leans in for a conspiratorial reveal. & How do gestures reinforce, contradict, or sharpen the spoken words? During split-screen listening, does the host's face (nodding vs. smirking/eye-rolling) cue the audience on whom to trust? How are bodily visuals used as evidence? \\ \hline
\end{longtable}

\subsection*{Quick Multimodal Transcription Notation Key (use inside transcript excerpts)}

\begin{longtable}{|p{0.15\textwidth}|p{0.15\textwidth}|p{0.15\textwidth}|p{0.15\textwidth}|p{0.15\textwidth}|p{0.15\textwidth}|}
\hline
Symbol & Meaning & Example in transcript & Symbol & Meaning & Example in transcript \\ \hline \endhead
\texttt{(.)} & Brief micro-pause (< 0.5 sec) & \texttt{woh (.) suparavijan lacking tha} & \texttt{$\uparrow$ / $\downarrow$} & Marked rising / falling pitch & \texttt{kyon$\uparrow$ FIR nahin karaee$\downarrow$} \\ \hline
\texttt{(1.5)} & Timed pause in seconds & \texttt{kisne (1.5) 'loota'?} & \texttt{WORD / \_word\_} & Loud volume / emphatic stress & \texttt{ye PAAP hai / \_bilkul\_ galat} \\ \hline
\texttt{---} & Abrupt stop, cut-off, or self-repair & \texttt{police ki taraf se--- abhi jaanch...} & \texttt{$^\circ$word$^\circ$} & Softened / whispered delivery & \texttt{$^\circ$ander ki baat ye hai$^\circ$} \\ \hline
\texttt{...} & Trailing off / unfinished sentence & \texttt{aap jaante hain ki kaun log...} & \texttt{[sigh] [cough] [scoff]} & Non-lexical vocalisations & \texttt{[sigh] 500 salon ke baad...} \\ \hline
\texttt{[overlap]} & Simultaneous / overlapping speech & \texttt{[Anjana: suniye---] [Guest: nahin!]} & \texttt{\{\{gesture: ...\}\}} & Hand gesture / facial cue / posture & \texttt{\{\{air-quotes on 'loota'\}\}} \\ \hline
\end{longtable}

\textbf{Methodological Reminder:} A non-verbal cue is not automatically manipulative or ideological---a cough, pause, or rapid cadence may reflect studio audio delay, live time pressure, or a speaker's habitual style. Always check (1) whether the cue clusters with verbal features, and (2) whether the speaker's delivery shifts noticeably across different actors or events (within-outlet baseline).\par\smallskip

\bigskip\hrule\bigskip

\subsection*{Event 1. Ram Mandir donation controversy (Non-Verbal \& Paralinguistic Layer)}

\begin{itemize}
  \item \textbf{Focus:} How do prosody, pauses, trailing sentences, vocalisations (sighs, coughs, scoffs), and body language frame the controversy---as solemn devotional sorrow, administrative scrutiny, or partisan combat?
  \item \textbf{Outlet:} Republic TV / Republic Bharat (Prime-Time Bulletin \& Debate Segment) | \textbf{Date:} 14 June 2021
  \item \textbf{Video title:} Ram Mandir Land Deal Allegations: Truth vs Conspiracy / Champat Rai \& Trust Response | \href{https://youtu.be/0H-bruoEkaY}{https://youtu.be/0H-bruoEkaY} | \textbf{Timestamp(s):} \texttt{[00:00]}--\texttt{[09:45]}
  \item \textbf{Analyst:} Baani Solanki (\texttt{2025114016}) | \textbf{Language:} Hindi / English mix (transcribed in Roman script)
  \item \textbf{Format / Visual setup:} \texttt{[x]} TV Panel Debate \texttt{[x]} TV Bulletin \texttt{[ ]} Solo YouTube Explainer \texttt{[ ]} Interview / Discussion
  \item \textbf{Primary speaker(s) observed:} Republic Anchor (standing/seated in high-alert red-and-yellow breaking-news studio), Champat Rai \& Anil Mishra (Trust spokespersons), and BJP vs. SP/AAP panelists
\end{itemize}

\subsubsection*{Event-specific non-verbal questions (Event 1)}

\noindent\textbf{1.} \textbf{Vocal \& gestural marking of core labels ("scam", "loota", "paap", "gadabdi", "avyavastha"): When the controversy is named, what non-verbal cues accompany the label? (Look for physical air-quotes, vocal scare-quotes, dramatic pauses before the word, hushed/reverential drops in volume, or shouted stress.)}\par\smallskip
\begin{itemize}
  \item \textbf{Notes:}
Republic TV never uses \texttt{"scam"} or \texttt{"ghotala"} as a factual description of the land transaction; instead, the word \texttt{"ghotala"} is always accompanied by \textbf{mocking vocal scare-quotes, physical air-quotes, or shouted counter-labels (\texttt{"SHADYANTRA"}, \texttt{"PROPAGANDA"}, \texttt{"JHOOTH"})}:\par\smallskip
  \item At \texttt{[00:25]}, the anchor introduces the opposition's allegation with a \texttt{(1.0)} pause and exaggerated two-handed air-quotes: \textit{"Ram Mandir ke naam par (1.0) 'ghotala'$\uparrow$?! Ya phir Mandir nirmaan rokne ki \_SAZISH\_$\downarrow$!"} (\texttt{[00:29]}).
  \item The word \textit{"\_SAZISH\_"} (conspiracy) receives maximum volume stress (\texttt{WORD}), a sharp downward fist-thump on the desk \texttt{\{\{gesture: right fist thumping desk\}\}}, and a fierce scowl, immediately inverting the story from a financial audit into an ideological war against the Temple.
\end{itemize}

\noindent\textbf{2.} \textbf{Differential demeanour toward actors (Temple, Trust leadership, counting clerks, donors, opposition accusers): How do the speaker's voice pitch, facial expression, and posture shift when referring to Prabhu Shri Ram / devotees vs. top figures (Champat Rai, Nripendra Misra, CM Yogi) vs. arrested counting staff vs. opposition leaders (Akhilesh Yadav, Awadhesh Prasad)?}\par\smallskip
\begin{itemize}
  \item \textbf{Notes:}
Republic exhibits extreme non-verbal polarisation between two moral camps:\par\smallskip
  \item \textbf{Prabhu Shri Ram, Devotees \& Trust Leadership (Champat Rai, Anil Mishra):} Treated as a single, indivisible sacred entity. When introducing Champat Rai's clarification (\texttt{[02:15]}), the anchor's pitch lowers into an \textbf{earnest, respectful, affirmative register}, posture upright, nodding continuously as Trust documents are shown (\texttt{"Trust ne ek-ek paise ka hisaab diya hai$\downarrow$"}, \texttt{[02:22]}).
  \item \textbf{Opposition Accusers (Sanjay Singh, Pawan Pandey, Akhilesh Yadav, Congress):} Addressed with \textbf{high-decibel combative shouting, mocking \texttt{[scoff]} vocalisations, rapid-fire finger-wagging directly at the lens, and aggressive forward leaning} (\texttt{"Jo log kal tak Ram ke astitva par sawaal uthate the, aaj unhe chande ki chinta ho rahi hai$\uparrow$?!"}, \texttt{[03:40]}).
\end{itemize}

\noindent\textbf{3.} \textbf{Solemn devotional grief vs. high-decibel partisan spectacle: How do pauses, sighs, speech rate, and hand gestures construct the segment's emotional register? Where is delivery slow, hushed, and mournful (sighs, bowed head, folded hands), and where does it escalate into rapid-fire shouting, desk-thumping, or mocking irony?}\par\smallskip
\begin{itemize}
  \item \textbf{Notes:}
Unlike Aaj Tak (which opens with slow, mournful devotional grief), Republic TV operates almost entirely in \textbf{high-decibel partisan spectacle} (\textasciitilde{}170--185 words per minute) right from \texttt{[00:00]}. There are virtually no quiet, mournful silences; pauses are short, sharp rhetorical gasps \texttt{[gasp]} used to reload before the next shouted clause (\texttt{[01:10]}, \texttt{[04:15]}). Hand gestures are large, percussive, and militaristic: continuous \textbf{desk-thumping, karate-chops in the air, and stabbing index-finger points} at the camera and at opposition split-screen boxes (\texttt{[05:20]}--\texttt{[08:30]}).\par\smallskip
\end{itemize}

\noindent\textbf{4.} \textbf{Non-verbal evaluation of financial scrutiny: When questions about unreceipted offerings, lack of frisking, or land deals are raised, how does the host/speaker react non-verbally (earnest nodding, leaning into the camera, open-palm appeals vs. eye-rolls, scoffs, dismissive hand waves, or incredulous rising pitch $\uparrow$)?}\par\smallskip
\begin{itemize}
  \item \textbf{Notes:}
Any question about the 5-minute price jump from Rs 2 crore to Rs 18.5 crore is met non-verbally with \textbf{dismissive hand waves, eye-rolls, loud \texttt{[scoff]} vocalisations, and incredulous rising pitch \texttt{$\uparrow$}}:\par\smallskip
  \item When an opposition panelist attempts to read out the \texttt{7:10 PM} vs. \texttt{7:15 PM} registry timestamps (\texttt{[05:45]}), the anchor rolls their eyes, waves the back of their hand dismissively across the desk \texttt{\{\{gesture: backhanded dismissive wave + eye-roll\}\}}, and shouts over them \texttt{[overlap]}: \textit{"Purana registered agreement aapko dikhai nahin deta$\uparrow$?! Ya padhna nahin aata$\uparrow$?!"} (\texttt{[05:52]}).
  \item Conversely, when the Trust spokesperson explains that the price was below market rate and paid digitally via RTGS (\texttt{[02:40]}--\texttt{[03:15]}), the anchor leans in with \textbf{earnest, vigorous head nods} \texttt{\{\{gesture: vigorous affirmative nodding\}\}} and taps the desk to endorse the point.
\end{itemize}

\noindent\textbf{5.} \textbf{Performing certainty vs. hedging and innuendo (pauses, stops, trailing off): Are allegations delivered with definitive falling intonation $\downarrow$, finger-counting, and pointing at on-screen evidence (CCTV, documents), or hedged with hesitant pauses (.), abrupt self-interruptions (---), throat-clearing [cough], and trailing sentences (...) that imply higher-level complicity without naming it?}\par\smallskip
\begin{itemize}
  \item \textbf{Notes:}
Republic performs \textbf{100\% epistemic certainty in exonerating the Trust} and \textbf{conspiratorial certainty in indicting the Opposition}:\par\smallskip
  \item \textbf{Exoneration of the Trust:} Delivered with booming falling intonation \texttt{$\downarrow$} and finger-counting (\texttt{"Number ek--- saara paisa bank se gaya$\downarrow$! Number do--- purana agreement 2011 se chal raha tha$\downarrow$!"}, \texttt{[03:05]}).
  \item \textbf{Innuendo / Conspiracy against Opposition:} Delivered via dramatic trailing clauses (\texttt{...}) and lowered, menacing pitch: \textit{"Chunav se theek pehle (.) Ayodhya ko badnaam karne ke liye ye khel khela ja raha hai... aur iske peechhe kaun si taqatein hain..."} (\texttt{[04:50]}), trailing off with a slow, menacing point at the lens \texttt{\{\{gesture: sustained finger-point at lens + narrowed eyes\}\}} to imply an "anti-Hindu / anti-national" cabal.
\end{itemize}

\noindent\textbf{6.} \textbf{Turn-taking, interruptions, and acoustic treatment of denials: In debates/interviews: who receives calm, uninterrupted speaking time and sympathetic nods, and who is cut off mid-sentence (---), talked over, or muted? In solo videos: how are clips of official denials framed via reaction pauses, raised eyebrows, or sarcastic chuckles?}\par\smallskip
\begin{itemize}
  \item \textbf{Notes:}
Turn-taking is \textbf{radically asymmetric and weaponised}:\par\smallskip
  \item \textbf{Trust Spokespersons \& BJP Defenders (\texttt{[02:15]}--\texttt{[03:30]}, \texttt{[06:30]}):} Receive 60--90 seconds of completely uninterrupted speaking time, clean solo audio, full-screen or dominant window placement, and affirmative verbal backchannels (\textit{"Bilkul sahi"}, \textit{"Yes"}).
  \item \textbf{Opposition Panelists (SP / AAP, \texttt{[05:35]}--\texttt{[08:15]}):} Cut off mid-sentence \texttt{---} within 5--8 seconds of speaking. The anchor repeatedly talks over them \texttt{[overlap]} at superior volume (\textit{"[Anchor: Jhooth mat boliye! Ek minute--- jhooth mat boliye!] [SP Guest: Kagaz toh dekhiye---]"}, \texttt{[06:05]}), while the studio audio engineer ducks the opposition guest's microphone volume by \textasciitilde{}10 dB so their rebuttal is reduced to unintelligible background mumbles.
\end{itemize}

\noindent\textbf{7.} \textbf{Performing the boundary between religious devotion and institutional critique: How does the presenter use vocal tone, sighs, and body language to signal personal reverence for the Temple (or secular objectivity) while criticising---or defending---the Trust's financial management?}\par\smallskip
\begin{itemize}
  \item \textbf{Notes:}
Republic \textbf{erases the boundary between religious devotion and institutional management} entirely. Through non-verbal fusion---placing a hand over the heart when saying \textit{"Shri Ram Janmabhoomi Teerth Kshetra Trust"} (\texttt{[01:35]}) and immediately pointing an angry finger when saying \textit{"Mandir par hamla"} (\texttt{[01:42]})---the anchor equates questioning Champat Rai's real-estate signatory decisions with attacking Lord Ram Himself. Any institutional audit is performed non-verbally as sacrilege.\par\smallskip
\end{itemize}

\noindent\textbf{8.} \textbf{Self-presentation as the "guardian of public trust": How does the speaker use direct eye-gaze, posture, moral sighs, or emphatic gestures (pointing at the viewer, hands on heart/desk) to position themselves as the authentic protector of devotees or public accountability?}\par\smallskip
\begin{itemize}
  \item \textbf{Notes:}
The anchor presents themselves as a \textbf{belligerent civilisational warrior defending Hindu devotees (\texttt{"karodon Ram bhakton ki aastha"}) against a hostile political cabal}. Standing/leaning aggressively into the camera, gripping a sheaf of papers in one hand and slapping them on the desk \texttt{\{\{gesture: slapping sheaf of papers on desk\}\}} (\texttt{[07:10]}), the host stares directly into the lens and shouts: \textit{"Desh ki janta is saazish ko \_kabhi\_ maaf nahin karegi$\downarrow$!"} (\texttt{[08:55]}).\par\smallskip
\end{itemize}

\subsubsection*{Apply the five non-verbal features (Event 1)}

\begin{longtable}{|p{0.31\textwidth}|p{0.31\textwidth}|p{0.31\textwidth}|}
\hline
Non-verbal feature & Timestamp + spoken anchor words + non-verbal notation & What it does in context (stance, framing, implication) \\ \hline \endhead
\textbf{1. Pauses, silences \& stops} & \texttt{[00:25]} \textit{"Ram Mandir ke naam par (1.0) 'ghotala'$\uparrow$?!"} / \texttt{[06:05]} \texttt{[overlap]} \textit{"[Anchor: Jhooth mat boliye--- chup rahiye!] [Guest: Do crore ka---]"} (guest mic faded) & The \texttt{(1.0)} pause sets up sarcastic air-quotes on \texttt{'ghotala'}, while forced silence via \texttt{[overlap]} shouting and audio ducking blocks the guest from stating the Rs 2 crore figure. \\ \hline
\textbf{2. Trailing sentences (\texttt{...})} & \texttt{[04:50]} \textit{"Chunav se theek pehle Ayodhya ko badnaam karne ke liye... aap jaante hain kaun log baukhlaye hue hain..."} & Trails off with a menacing stare to presuppose that anyone questioning the land deal is driven by electoral panic and anti-temple malice. \\ \hline
\textbf{3. Intonation, pitch \& volume} & \texttt{[00:29]} \textit{"Mandir nirmaan rokne ki \_SAZISH\_$\downarrow$!"} (maximum volume \texttt{WORD} + falling finality \texttt{$\downarrow$}) vs. \texttt{[05:52]} \textit{"Padhna nahin aata$\uparrow$?!"} (sharp mocking rise \texttt{$\uparrow$}) & Uses booming falling stress \texttt{$\downarrow$} to declare an opposition conspiracy as established fact, and mocking rising pitch \texttt{$\uparrow$} to humiliate dissenting panelists. \\ \hline
\textbf{4. Coughs, sighs \& vocalisations} & \texttt{[01:10]} \texttt{[[gasp]]} (sharp intake of breath between rapid shouted clauses) / \texttt{[03:45]} \texttt{[[scoff]]} \textit{"Aaj inhe chande ki chinta ho rahi hai!"} & Rapid \texttt{[gasp]} intakes sustain high-arousal urgency; the loud \texttt{[scoff]} ridicules the sincerity of opposition leaders' concern for donors. \\ \hline
\textbf{5. Hand gestures, face \& posture} & \texttt{[00:25]} \texttt{\{\{gesture: two-handed air-quotes on 'ghotala'\}\}} / \texttt{[00:29]} \texttt{\{\{gesture: right fist thumping desk\}\}} / \texttt{[07:10]} \texttt{\{\{gesture: slapping papers on desk + pointing at lens\}\}} & Air-quotes delegitimise the allegation; desk-thumping and paper-slapping perform prosecutorial dominance and righteous nationalist anger. \\ \hline
\end{longtable}

\subsubsection*{Short summary (Event 1 Non-Verbal Synthesis)}

\begin{itemize}
  \item \textbf{What overall interpretation or emotional stance does the non-verbal delivery encourage?}
The non-verbal delivery encourages \textbf{combative, defensive outrage against the Opposition}. By pairing air-quotes on \textit{"'ghotala'"} with fist-thumping stress on \textit{"\_SAZISH\_"} and relentless \texttt{[overlap]} interruptions of critics, the broadcast leaves zero room for administrative doubt, framing the entire controversy as a fabricated attack on the Ram Mandir.\par\smallskip
  \item \textbf{Whose voice or physical presence is acoustically/visually foregrounded, undermined, or silenced?}
\textbf{The studio anchor and Trust leadership (Champat Rai / Anil Mishra)} are acoustically and visually foregrounded with uninterrupted solo segments and vigorous affirmative nods. \textbf{Opposition spokespersons (SP/AAP)} are systematically \textbf{silenced}---cut off mid-sentence \texttt{---} within seconds, shouted over \texttt{[overlap]}, and acoustically ducked in the studio mix (\texttt{[06:05]}).\par\smallskip
  \item \textbf{What evidence supports your non-verbal reading? What complicates it (e.g., speaker baseline, format)?}
The turn-taking data---60+ seconds of uninterrupted calm audio for Trust defenders (\texttt{[02:15]}--\texttt{[03:30]}) versus <8 seconds before \texttt{[overlap]} interruption for critics (\texttt{[05:35]}--\texttt{[08:15]})---is objective acoustic proof of gatekeeping. While high-decibel shouting (\textasciitilde{}175 WPM) is Republic TV's general prime-time baseline, the complete absence of interruptions when Trust representatives speak shows that the anchor's aggression is selectively deployed.\par\smallskip
\end{itemize}

\bigskip\hrule\bigskip

\subsection*{Event 2. Students' protest regarding NEET paper leak (Non-Verbal \& Paralinguistic Layer)}

\begin{itemize}
  \item \textbf{Focus:} How do prosody, pauses, trailing sentences, vocalisations, and visual body language construct the legitimacy of student protesters, police use of force, and state accountability?
  \item \textbf{Outlet:} Republic TV (\textit{The Debate with Arnab Goswami}) | \textbf{Date:} June 2024
  \item \textbf{Video title:} NEET Paper Leak \& Student Protest Debate | Arnab Goswami on NTA \& Political Row | \href{https://youtu.be/HTvb6E\_1Qdc}{https://youtu.be/HTvb6E\_1Qdc} | \textbf{Timestamp(s):} \texttt{[00:00]}--\texttt{[36:20]}
  \item \textbf{Analyst:} Baani Solanki (\texttt{2025114016}) | \textbf{Language:} English (with occasional Hindi terms in Roman script)
  \item \textbf{Format / Visual setup:} \texttt{[x]} TV Panel Debate \texttt{[ ]} TV Bulletin \texttt{[ ]} Solo YouTube Explainer \texttt{[ ]} Interview / Discussion
  \item \textbf{Primary speaker(s) observed:} Arnab Goswami (studio anchor seated/leaning forward at desk), surrounded by an 8-window split-screen grid of Student Representatives, Educators, BJP/NDA Spokespersons, and Opposition (Congress/RJD/SP) Panelists
\end{itemize}

\subsubsection*{Event-specific non-verbal questions (Event 2)}

\noindent\textbf{1.} \textbf{Prosodic \& gestural marking of participant labels ("students/bachche" vs. "protesters" vs. "upadravi/goons"): How do vocal tone, volume, facial expression, and hand gestures shift across labels? (e.g., softened, empathetic pitch for "students" vs. harsh, staccato stress, scowls, or accusatory finger-pointing on "upadraviyon", "goons", or "anti-social elements".)}\par\smallskip
\begin{itemize}
  \item \textbf{Notes:}
Arnab Goswami creates a dramatic \textbf{three-way prosodic and gestural split} across labels:\par\smallskip
  \item \textbf{Genuine NEET Aspirants (\texttt{"young students"}, \texttt{"24 lakh children"}):} Delivered with a \textbf{deliberately slowed, earnest, protective cadence} (\texttt{[00:35]}, \texttt{[04:10]}), both hands held open toward the camera \texttt{\{\{gesture: both palms open and extended toward lens\}\}} and a deeply furrowed brow of paternal concern.
  \item \textbf{Bureaucrats (\texttt{"NTA babus"}, \texttt{"faceless bureaucrats"}):} Delivered with biting, staccato contempt (\texttt{[02:15]}), jabbing his pen downward into the desk \texttt{\{\{gesture: repeated pen-jabs onto desk surface\}\}}.
  \item \textbf{Opposition Youth-Wing Protesters \& Politicians (\texttt{"political vultures"}, \texttt{"hijackers"}, \texttt{"RJD-Congress lobby"}):} Delivered at \textbf{maximum shouted volume (\texttt{WORD}), rapid-fire pace (\textasciitilde{}190 WPM), red-faced scowls, and aggressive two-handed finger-pointing} (\texttt{"Stop hijacking the students' pain! You are \_POLITICAL VULTURES\_$\downarrow$!"}, \texttt{[14:25]}).
\end{itemize}

\noindent\textbf{2.} \textbf{Delivery of students' stated demands: When the protesters' core demands (Education Minister's resignation, compensation, release of Sonam Wangchuk, exam reform) are mentioned, are they delivered slowly with deliberate pauses and finger-counting, rushed through in a flat monotone, mocked with air-quotes/scoffs, or trailed off (...)?}\par\smallskip
\begin{itemize}
  \item \textbf{Notes:}
Arnab selectively \textbf{splits the students' demands into "acceptable administrative demands" vs. "unacceptable political demands"}:\par\smallskip
  \item \textbf{Administrative Demands (Scrapping grace marks, firing NTA officials, CBI probe, \texttt{[03:10]}--\texttt{[04:45]}):} Delivered \textbf{slowly and forcefully with finger-counting} (\texttt{"Point one (.) scrap the grace marks$\downarrow$! Point two (.) overhaul the NTA from top to bottom$\downarrow$!"}, \texttt{[03:22]}), thumping his fist on each number.
  \item \textbf{Political Demands (Resignation of Education Minister Dharmendra Pradhan or street blockades by youth wings, \texttt{[15:10]}):} Immediately \textbf{mocked with a loud \texttt{[scoff]}, dismissive hand-waving, and shouted interruptions} (\texttt{"You don't care about NEET! You only want a political resignation for your election narrative!"}, \texttt{[15:28]}).
\end{itemize}

\noindent\textbf{3.} \textbf{Where does vocal \& kinetic energy peak---grievance, street disruption, or electoral horse-race?: Track where the speaker's volume, pitch range, pace, and gesture intensity reach their climax: during accounts of the paper leak/student suicides, during street clashes and barricade breaches, or during partisan BJP-vs-Opposition sparring.}\par\smallskip
\begin{itemize}
  \item \textbf{Notes:}
While Arnab opens with high vocal energy against the NTA (\texttt{[00:00]}--\texttt{[06:00]}), the \textbf{absolute acoustic and kinetic climax of the 36-minute broadcast occurs during partisan BJP-vs-Opposition sparring over the Bihar paper-leak nexus (\texttt{[12:00]}--\texttt{[28:00]})}:\par\smallskip
  \item When linking the arrested Bihar mastermind Sikandar Yadavendu to RJD leader Tejashwi Yadav's personal secretary Pritam Kumar (\texttt{[13:40]}--\texttt{[18:50]}), Arnab's volume hits peak saturation (\texttt{WORD}), he rises halfway out of his chair \texttt{\{\{gesture: rising half-standing from chair, leaning over desk\}\}}, waves printed call-detail/guest-house papers at the camera, and shouts non-stop for 4 minutes at the RJD/Congress panelists.
\end{itemize}

\noindent\textbf{4.} \textbf{Non-verbal attribution of responsibility (sighs, stops, head shakes, pointing): When examination authorities (NTA), the Education Minister, police, opposition politicians, or students are named as responsible, what non-verbal cues accompany the attribution (exasperated sighs, mocking laughs, abrupt self-stops ---, or direct finger-pointing at the lens)?}\par\smallskip
\begin{itemize}
  \item \textbf{Notes:}
Arnab performs a classic \textbf{non-verbal blame-deflection manoeuvre} that shields the Union Government while attacking both the NTA bureaucracy and the Opposition:\par\smallskip
  \item \textbf{NTA Bureaucrats:} Blamed with exasperated head shakes, \texttt{[sigh]} vocalisations, and pen-pointing (\texttt{"NTA has become a law unto itself"}, \texttt{[02:45]}), isolating failure to an "autonomous agency".
  \item \textbf{Union Education Minister / Modi Government:} Whenever an opposition panelist names Dharmendra Pradhan or PM Modi (\texttt{[16:10]}), Arnab abruptly stops them \texttt{---}, shakes his head vigorously, and pivots (\texttt{"The government has already ordered a CBI probe and enacted a strict anti-paper-leak law$\downarrow$"}, \texttt{[16:22]}).
  \item \textbf{Opposition Politicians (Tejashwi Yadav / RJD):} Blamed for the actual Patna paper leak with \textbf{direct finger-pointing at the lens, mocking \texttt{[chuckle]}/\texttt{[scoff]} vocalisations, and desk-slapping} (\texttt{"Who booked the NHAI guest house in Patna$\uparrow$?! \_PRITAM KUMAR\_$\downarrow$!"}, \texttt{[14:05]}).
\end{itemize}

\noindent\textbf{5.} \textbf{Prosody of questions---open inquiry vs. loaded rhetorical challenge: Are questions about the protest delivered with level, open intonation, or with sharp rising pitch $\uparrow$, sarcastic stress, raised eyebrows, and trailing clauses (...) that presuppose protesters are violent or politically manipulated?}\par\smallskip
\begin{itemize}
  \item \textbf{Notes:}
  \item \textbf{To Student / Educator Guests (\texttt{[07:15]}, \texttt{[09:30]}):} Arnab switches to a \textbf{remarkably soft, solicitous, open-inquiry prosody} (\texttt{"Tell me, young man, what did you feel when you saw 67 toppers$\uparrow$?"}, \texttt{[07:20]}), leaning his chin on his hand and nodding empathetically \texttt{\{\{gesture: chin resting on hand + slow empathetic nod\}\}}.
  \item \textbf{To Opposition / Youth-Wing Protesters (\texttt{[13:50]}, \texttt{[19:20]}):} Delivered as \textbf{explosive, loaded rhetorical challenges} with sharp rising pitch \texttt{$\uparrow$}, sarcastic stress, and accusatory trailing off (\texttt{...}): \textit{"Why is Rahul Gandhi silent on the Bihar guest house link$\uparrow$?! Because the trail leads straight to..."} (\texttt{[14:48]}).
\end{itemize}

\noindent\textbf{6.} \textbf{Visual \& paralinguistic performance of injury and violence: How are claims of violence substantiated non-verbally on screen? (Compare camera close-ups and pained vocal tone on police bandages/fractured hands vs. visual/acoustic treatment of student head injuries, lathi-charge footage, or defensive stammers/pauses when an officer or guest is asked who used force first.)}\par\smallskip
\begin{itemize}
  \item \textbf{Notes:}
Republic splits visual and paralinguistic treatment between \textbf{"innocent aspirants in studios"} and \textbf{"rowdy political cadres on streets"}:\par\smallskip
  \item Ordinary NEET candidates in the studio grid receive steady close-up frames and pained vocal sympathy from Arnab (\texttt{[07:15]}--\texttt{[10:30]}).
  \item Street protest clashes (\texttt{[20:10]}) are framed via split-screen loops showing Congress/NSUI/Left flags at barricades while Arnab points at the party flags on screen \texttt{\{\{gesture: pointing at party flags in B-roll window\}\}} and shouts: \textit{"Look at the flags! These are \_NOT\_ students, these are Congress workers creating anarchy on the streets$\downarrow$!"} (\texttt{[20:25]}), completely dismissing any claims of police excess against youth-wing marchers.
\end{itemize}

\noindent\textbf{7.} \textbf{Asymmetric delivery of police actions vs. protester actions: Compare the vocal cadence, pauses, and hand gestures used when describing police actions (e.g., calm, matter-of-fact administrative tone on "tear gas chhode gae" / "preventive arrangement") versus protester actions (alarmist pitch, striking gestures, or ironic softening as "retaliation").}\par\smallskip
\begin{itemize}
  \item \textbf{Notes:}
  \item \textbf{Police \& CBI Actions:} Narrated in a \textbf{resolute, authoritative, approving cadence} with firm downward hand-chops (\texttt{"Bihar Police and EOU have cracked the nexus$\downarrow$! CBI is taking over$\downarrow$!"}, \texttt{[13:15]}), portraying law enforcement as heroic investigators uncovering opposition-linked mafia rings.
  \item \textbf{Street Protester Actions (Opposition Youth Wings):} Narrated with \textbf{alarmist pitch \texttt{$\uparrow$}, mocking \texttt{[scoff]} vocalisations, and striking gestures} (\texttt{"Pelting stones and climbing barricades will not solve NEET---it is cheap political drama!"}, \texttt{[20:40]}).
\end{itemize}

\noindent\textbf{8.} \textbf{Acoustic \& visual space granted to students' own voices: Are protesting students heard in their own voice with clear audio, steady close-ups, and empathetic listening nods, or are they reduced to noisy background chaos B-roll while studio anchors/politicians speak over them?}\par\smallskip
\begin{itemize}
  \item \textbf{Notes:}
Arnab grants \textbf{clean, protected acoustic space to non-political NEET students/educators} in the first third of the show (\texttt{[07:10]}--\texttt{[11:45]}), explicitly shouting at politicians to stay silent while the student speaks (\textit{"Quiet! Let the student speak! Don't interrupt the student!"}, \texttt{[08:05]}, holding up a flat palm like a traffic stop sign \texttt{\{\{gesture: right palm held up vertically to silence politicians\}\}}). However, once the debate shifts to the street protests and the Bihar RJD link (\texttt{[12:00]}--\texttt{[36:20]}), students in the 8-box grid are \textbf{visually sidelined and acoustically drowned out} for over 20 minutes by non-stop \texttt{[overlap]} shouting between Arnab, BJP, and RJD/Congress panelists.\par\smallskip
\end{itemize}

\noindent\textbf{9.} \textbf{Non-verbal shift between the NEET paper-leak issue and the street protest: Does the speaker's non-verbal demeanour shift from sympathetic concern (sombre tone, sighs, solidarity gestures) when discussing the exam leak itself to combative, authoritarian, or dismissive delivery once students march toward Parliament?}\par\smallskip
\begin{itemize}
  \item \textbf{Notes:}
\textbf{Yes---an extreme bimodal shift occurs.} During the exam-anomaly segment (\texttt{[00:00]}--\texttt{[11:45]}), Arnab performs \textbf{paternal solidarity} with students and anger at NTA bureaucrats (leaning in softly to listen to students, furrowing his brow). The moment the topic shifts to the street protests and political accountability (\texttt{[12:00]} onward), he pivots into \textbf{hyper-combative partisan warfare}: shouting at \textasciitilde{}190 WPM, slapping papers on the desk, and accusing street protesters of being Congress/RJD puppets trying to destabilise the government.\par\smallskip
\end{itemize}

\subsubsection*{Apply the five non-verbal features (Event 2)}

\begin{longtable}{|p{0.31\textwidth}|p{0.31\textwidth}|p{0.31\textwidth}|}
\hline
Non-verbal feature & Timestamp + spoken anchor words + non-verbal notation & What it does in context (stance, framing, implication) \\ \hline \endhead
\textbf{1. Pauses, silences \& stops} & \texttt{[08:05]} \textit{"Quiet! (1.0) Let the student speak$\downarrow$!"} / \texttt{[14:02]} \textit{"Who booked the NHAI guest house in Patna? (1.5) \_PRITAM KUMAR\_, PA to Tejashwi Yadav$\downarrow$!"} & The \texttt{(1.0)} command silence positions Arnab as the protector of students; the \texttt{(1.5)} dramatic pause before naming Tejashwi's aide turns the NEET debate into an indictment of the Opposition. \\ \hline
\textbf{2. Trailing sentences (\texttt{...})} & \texttt{[14:48]} \textit{"Why is Rahul Gandhi suddenly silent on the Patna arrest$\uparrow$?! Because the paper-leak trail leads straight to his alliance partner..."} & Trails off while pointing at the lens, inviting the viewer to conclude that the Opposition itself engineered the paper leak. \\ \hline
\textbf{3. Intonation, pitch \& volume} & \texttt{[07:20]} \textit{"$^\circ$Tell me, young man, how hard did you prepare$^\circ$?"} vs. \texttt{[14:25]} \textit{"You are \_POLITICAL VULTURES\_ hijacking the students!"} & Extreme prosodic contrast between hushed, gentle pitch \texttt{$^\circ$...$^\circ$} for individual students and maximum shouted volume \texttt{WORD} against opposition protesters. \\ \hline
\textbf{4. Coughs, sighs \& vocalisations} & \texttt{[02:45]} \texttt{[[sigh]]} \textit{"My heart goes out to the \texttt{24 lakh} children..."} / \texttt{[15:25]} \texttt{[[scoff] + [chuckle]]} (at demand for Education Minister's resignation) & The opening \texttt{[sigh]} performs populist empathy; the mocking \texttt{[scoff]} dismisses calls for ministerial accountability as partisan opportunism. \\ \hline
\textbf{5. Hand gestures, face \& posture} & \texttt{[08:05]} \texttt{\{\{gesture: vertical stop-palm toward politicians, open palm toward student\}\}} / \texttt{[14:10]} \texttt{\{\{gesture: rising half-standing, waving Bihar guest-house papers\}\}} & The two-handed traffic-cop gesture visually separates "pure students" from "dirty politicians"; waving the Bihar papers redirects anger away from the Centre toward the RJD. \\ \hline
\end{longtable}

\subsubsection*{Short summary (Event 2 Non-Verbal Synthesis)}

\begin{itemize}
  \item \textbf{What overall interpretation or emotional stance does the non-verbal delivery encourage?}
The non-verbal delivery encourages \textbf{populist sympathy for individual exam-takers, anger at nameless "NTA bureaucrats", and fierce partisan hostility toward street protesters and the Opposition}. By splitting "innocent students" from "political hijackers", Arnab channels public anger over the NEET scam away from the Union Government and onto the Bihar Opposition (\texttt{RJD/Congress}).\par\smallskip
  \item \textbf{Whose voice or physical presence is acoustically/visually foregrounded, undermined, or silenced?}
\textbf{Arnab Goswami} dominates the acoustic floor throughout. \textbf{Individual non-political students} are foregrounded briefly (\texttt{[07:10]}--\texttt{[11:45]}) with clean audio and empathetic nods, after which they sit silently in small grid boxes while \textbf{Opposition spokespersons and street protesters} are shouted over \texttt{[overlap]}, mocked with scoffs, and visually tagged with party flags.\par\smallskip
  \item \textbf{What evidence supports your non-verbal reading? What complicates it (e.g., speaker baseline, format)?}
The stark acoustic contrast between \texttt{[07:20]} (\textasciitilde{}115 WPM, softened volume \texttt{$^\circ$...$^\circ$}, chin-on-hand listening posture for a student) and \texttt{[14:05]} (\textasciitilde{}190 WPM, shouted \texttt{WORD}, half-standing desk-slapping on the RJD guest-house link) proves a strategic within-broadcast shift. While Arnab's shouting baseline is famous across all \textit{Debate} episodes, his ability to drop into quiet, uninterrupted listening for a student guest shows that his high-decibel \texttt{[overlap]} interruptions of opposition guests are deliberate turn-taking control.\par\smallskip
  \item \textbf{How do the non-verbal cues interact with the spoken words? (Do they reinforce, dramatise, soften/hedge, invert meaning via sarcasm/irony, or leave an accusation unspoken via trailing off?)}
His non-verbal cues \textbf{dramatise bureaucratic failure} (pen-jabbing on \textit{"NTA babus"}), \textbf{soften central government culpability} (calm, reassuring nods when mentioning the CBI probe at \texttt{[16:22]}), \textbf{invert street protests via visual pointing} (pointing at party flags to rebrand protesters as \textit{"Congress cadres"} at \texttt{[20:25]}), and \textbf{implicate the Opposition via dramatic \texttt{(1.5)} pauses and trailing off} (\texttt{...} at \texttt{[14:48]}).\par\smallskip
  \item \textbf{What context, speaker baseline, or format factors support or complicate this interpretation? (e.g., Is this the speaker's normal style across all videos, or a marked shift on this specific actor/claim?)}
Comparing Event 1 and Event 2 shows Republic TV's signature \textbf{"populist-nationalist shield"}: in Event 1 (where a sacred institution is accused), the channel denies any wrongdoing whatsoever (\texttt{"100\% sazish"}). In Event 2 (where \texttt{24 lakh} students and middle-class families are angry and denial is impossible), Arnab adopts a softer, empathetic posture toward the victims for the first 10 minutes before deploying his standard high-decibel prosecutorial arsenal to blame autonomous bureaucrats and an Opposition-linked paper-leak nexus in Bihar.\par\smallskip
  \item \textbf{How does another outlet perform the same moment, actor, or claim non-verbally?}
Whereas \textbf{Dhruv Rathee} (\texttt{Event 2}) and \textbf{The Wire} (\texttt{Event 2}) use document-pointing and rising pitch \texttt{$\uparrow$} to connect the paper leaks in Bihar and Gujarat (Godhra) directly to the Union Education Minister and ruling-party states, \textbf{Republic TV} (\texttt{Event 2}) waves documents exclusively about the Patna NHAI guest house (\texttt{[14:05]}) with booming stress on \textit{"\_PRITAM KUMAR\_"} (Tejashwi Yadav's aide), using non-verbal spectacle to turn a central exam failure into an opposition scandal.\par\smallskip
  \item \textbf{My interpretation, with reasons:}
Republic TV's multimodal discourse across both events demonstrates \textbf{acoustic domination and strategic blame redirection}. Through \texttt{[overlap]} shouting, selective mic-ducking, desk-thumping, and split-screen visual framing, the channel neutralises threats to establishment legitimacy---either by branding financial audit questions as an "anti-Mandir conspiracy" (Event 1) or by splitting sympathetic exam aspirants from "vulture" street protesters and pinning the paper leak on the Opposition (Event 2).\par\smallskip
\end{itemize}

\end{document}
